\documentclass[11pt,a4paper]{article}
\usepackage{iftex}
\ifPDFTeX
  \usepackage[T1]{fontenc}
  \usepackage[utf8]{inputenc}
  \usepackage{lmodern}
\else
  \usepackage{fontspec}
  \setmainfont[BoldFont=lmroman10-bold.otf,ItalicFont=lmroman10-italic.otf,BoldItalicFont=lmroman10-bolditalic.otf]{lmroman10-regular.otf}
  \setmonofont{lmmono10-regular.otf}
\fi
\usepackage[polish]{babel}
\usepackage[margin=2.5cm]{geometry}
\usepackage{xcolor}
\usepackage{fancyvrb}
\usepackage{fvextra}
\usepackage{newunicodechar}
% Znaki spoza fontu Latin Modern (greka, indeksy, symbole) składane jako matematyka.
\newunicodechar{α}{\ensuremath{\alpha}}
\newunicodechar{γ}{\ensuremath{\gamma}}
\newunicodechar{ε}{\ensuremath{\varepsilon}}
\newunicodechar{σ}{\ensuremath{\sigma}}
\newunicodechar{₀}{\ensuremath{_0}}
\newunicodechar{≤}{\ensuremath{\leq}}
\newunicodechar{≥}{\ensuremath{\geq}}
\newunicodechar{×}{\ensuremath{\times}}
\newunicodechar{−}{\ensuremath{-}}
\newunicodechar{→}{\ensuremath{\rightarrow}}
\usepackage{hyperref}
\hypersetup{colorlinks=true,linkcolor=black,urlcolor=blue}
% Treść rozmów: zwykła czcionka, zawijanie wierszy bez symboli strzałek, ramka.
\fvset{fontfamily=rm,fontsize=\normalsize,breaklines=true,breakanywhere=true,
  breaksymbolleft={},breaksymbolright={},breakindent=0pt,
  frame=leftline,framerule=1.5pt,framesep=8pt,rulecolor=\color{gray!60},xleftmargin=4pt}
\setlength{\parindent}{0pt}
\setlength{\parskip}{0.4em}
\newcommand{\metawiersz}[2]{\textbf{#1:} #2\par}
\title{Dziennik wsparcia AI}
\author{Projekt: Uczenie się ze wzmocnieniem 2026/2027}
\date{}
\begin{document}
\maketitle
Dziennik zawiera rozmowy z narzędziami AI wykorzystanymi w projekcie
(zgodnie z P00, sekcja 9): datę, narzędzie i model, cel, zakres pomocy
oraz pełną treść promptu i odpowiedzi. Daty w strefie Europe/Warsaw.

\tableofcontents
\clearpage
% Generowane przez scripts/log_ai.py; nie edytuj ręcznie.
\section{2026-10-21 -- Formalizacja MDP}
\label{20261021-ca5c95bfa1fe}
\metawiersz{Data}{2026-10-21}
\metawiersz{Narzędzie}{Claude (Cowork) (claude-opus-5-5)}
\metawiersz{Cel}{Formalizacja MDP}
\metawiersz{Zakres pomocy}{Claude przedstawił opcje i zalecenia, student wybrał warianty; Claude zredagował sekcje raportu z formalizacją}
\metawiersz{Weryfikacja}{Decyzje projektowe podjął student. Weryfikacja treści raportu: do uzupełnienia.}
\subsection*{Prompt}
\VerbatimInput{entries/20261021-ca5c95bfa1fe/prompt.txt}
\subsection*{Odpowiedź}
\VerbatimInput{entries/20261021-ca5c95bfa1fe/response.txt}
\bigskip\hrule\bigskip

\section{2026-10-21 -- Plan eksperymentów i agenci Q-learning i SARSA}
\label{20261021-8169fa9df9dd}
\metawiersz{Data}{2026-10-21}
\metawiersz{Narzędzie}{Claude (Cowork) (claude-opus-5-5)}
\metawiersz{Cel}{Plan eksperymentów i agenci Q-learning i SARSA}
\metawiersz{Zakres pomocy}{Student wybrał algorytmy, liczbę ziaren i eksperyment dodatkowy; Claude napisał kod agentów, testy i sekcje raportu}
\metawiersz{Weryfikacja}{Wybory dokonane przez studenta. Weryfikacja kodu i raportu: do uzupełnienia.}
\subsection*{Prompt}
\VerbatimInput{entries/20261021-8169fa9df9dd/prompt.txt}
\subsection*{Odpowiedź}
\VerbatimInput{entries/20261021-8169fa9df9dd/response.txt}
\bigskip\hrule\bigskip

\section{2026-10-21 -- Dziennik i publikacja w repozytorium}
\label{20261021-ce1b1f57a2fb}
\metawiersz{Data}{2026-10-21}
\metawiersz{Narzędzie}{Claude (Cowork) (claude-opus-5-5)}
\metawiersz{Cel}{Dziennik i publikacja w repozytorium}
\metawiersz{Zakres pomocy}{Potwierdzenie wpisów w dzienniku; instrukcja commitu i pusha}
\metawiersz{Weryfikacja}{Do uzupełnienia.}
\subsection*{Prompt}
\VerbatimInput{entries/20261021-ce1b1f57a2fb/prompt.txt}
\subsection*{Odpowiedź}
\VerbatimInput{entries/20261021-ce1b1f57a2fb/response.txt}
\bigskip\hrule\bigskip

\section{2026-10-21 -- Publikacja w repozytorium GitHub}
\label{20261021-8a46b01167d0}
\metawiersz{Data}{2026-10-21}
\metawiersz{Narzędzie}{Claude (Cowork) (claude-opus-5-5)}
\metawiersz{Cel}{Publikacja w repozytorium GitHub}
\metawiersz{Zakres pomocy}{Kontrola zmian przed commitem; instrukcja commitu i pusha}
\metawiersz{Weryfikacja}{Do uzupełnienia.}
\subsection*{Prompt}
\VerbatimInput{entries/20261021-8a46b01167d0/prompt.txt}
\subsection*{Odpowiedź}
\VerbatimInput{entries/20261021-8a46b01167d0/response.txt}
\bigskip\hrule\bigskip


\end{document}
